% Source-derived chapter 07: Topology of Hitchin fibrations
% Original source: hitchin-fibrations-abelian-surfaces-pw
\chapter{Topology of Hitchin fibrations}
\label{hitchin-fibrations-abelian-surfaces-pw-chapter-07}
\source{hitchin-fibrations-abelian-surfaces-pw}

\begin{remark}[Source section]
\label{hitchin-fibrations-abelian-surfaces-pw-chapter-07-source}
\source{hitchin-fibrations-abelian-surfaces-pw}
This chapter records the source section; no Lean declaration has yet been checked for this material.
\notready
\end{remark}

% The source section title was: Topology of Hitchin fibrations
\subsection{Overview} Throughout the section, we assume $C$ is a nonsingular curve of genus $g$ embedded into an abelian surface $A$. We study a specialization morphism
\begin{equation}\label{sp}
\source{hitchin-fibrations-abelian-surfaces-pw}
    \mathrm{sp}^!: H^*(\CM_{\beta, A}, \BQ) \to H^*(\CM_{\mathrm{Dol}}, \BQ)
\end{equation}
where $\CM_{\mathrm{Dol}}$ is the moduli of stable Higgs bundles of rank $n$ and degree 1, and
\[
\beta = n[C] \in H_2(A, \BZ).
\]
Then we deduce Theorems \ref{genus_2}, \ref{P=W_even}, and \ref{P=W_odd} from Theorem \ref{taut_abelian} via the morphism (\ref{sp}).

\subsection{Deformation to the normal cone}\label{Section 4.2}
\source{hitchin-fibrations-abelian-surfaces-pw}

The moduli space $\CM_{\mathrm{Dol}}$ of rank $n$ stable Higgs bundles can be realized as the moduli space of $1$-dimensional Gieseker-stable sheaves $\CF$ on the surface $T^*C$ with
\[
 [\mathrm{supp}(\CF)] = n[C] \in H_2(T^*C, \BZ), \quad \chi(\CF) = 1;
\]
see \cite{BNR}. We describe $\CM_{\mathrm{Dol}}$ as the ``limit'' of a trivial family of $\CM_{\beta,A}$ as follows.

Assume $T = \BP^1$ and $T^\circ = \BP^1 \setminus 0 = \BC$. Let
\begin{equation}\label{normal_cone}
\source{hitchin-fibrations-abelian-surfaces-pw}
 \CS = \mathrm{Bl}_{C\times 0}(A \times T) \setminus (A\times 0) \to T
\end{equation}
be the total space of the deformation to the normal cone associated with the embedding $j_C:C\hookrightarrow A$. The central fiber of (\ref{normal_cone}) is
\[
\CS_0 = T^*C \to 0 \in T,
\]
and the restriction over $T^\circ$ is a trivial fibration
\[
A \times T^\circ \to  T^\circ \subset T.
\]


We associate to $\beta$ a family of homology classes
\[
\beta_t = n[C] \in H_2(\CS_t, \BZ).
\]
Let $\CM \rightarrow T$ be the (coarse) relative moduli space which parametrizes, for each $t \in T$, pure one-dimensional Gieseker-stable sheaves $\CF$ on $\CS_t$ such that $\CF$ has proper support in the class $\beta_t$ and $\chi(\CF)=1$. Similarly, let $\CB \rightarrow T$ denote the component of the relative Hilbert scheme which parametrizes Cartier divisors in $\CS_t$ with proper support in the class $\beta_t$.  Following \cite[Section 5.3]{GIT}, we have a proper morphism
\begin{equation*}
\text{[Diagram omitted; see source PDF.]}
\end{equation*}
over the base $T$, which, on the level of points, sends a sheaf $\CF$ on $\CS_t$ to its cycle-theoretic support, viewed as an element in $\CB_t$.
Clearly, both $\CM$ and $\CB$ become trivial families when restricted over $T^\circ$,
\begin{equation*}
\text{[Diagram omitted; see source PDF.]}
\end{equation*}
The fibers over $0 \in T$ recover the moduli space $\CM_{\mathrm{Dol}}$, the Hitchin base $\Lambda$, and the Hitchin fibration (\ref{Hitchin_fib}).

The following lemma is standard, but we give a brief proof due to lack of adequate reference.
\begin{lem}
\source{hitchin-fibrations-abelian-surfaces-pw}
\notready\label{lem4.1}
\begin{enumerate}
\item[(i)] Both $\CM$ and $\CB$ are irreducible varieties and smooth over $T$.
\item[(ii)] There exists a universal family $\CF_T$ of 1-dimensional Gieseker-stable sheaves on $\CM\times_{T} \CS$.
\end{enumerate}
\end{lem}
\begin{proof}

For (i), we first argue for $\CB$. For notational convenience, we define
\[
\tilde{g} = \mathrm{dim}(\CB_t) = n^2(g-1)+1, \quad \forall~t \in T.
\]
We observe that the closure of
\[
\CB^{\circ} = B \times T^\circ
\]
is irreducible of dimension $\tilde{g}+1$, so the intersection
\[
\overline{\CB^{\circ}}\cap \CB_0 \subset \CB_0
\]
must have at least dimension $\tilde{g}$ at every closed point.  This intersection is non-empty since it contains the divisor $n[C]$ which clearly deforms to the generic fiber. Since $\CB_0$ is irreducible of dimension $\tilde{g}$, we have $\CB_0 \subset \overline{\CB^{\circ}}$. Therefore $\CB$ is irreducible. This implies the flatness of $\CB$ over $T$. Furthermore, since its fibers are all nonsingular, the morphism $\CB \to T$ is smooth.

The same argument applies for $\CM$. The only thing to check is that there exists a point in the central fiber $\CM_0$ which deforms to the generic fiber.  By the last paragraph, we can choose a nonsingular curve $Z_0 \subset \CS_0$ which deforms to $Z_t \subset \CS_t$.  Any line bundle on $Z_0$ with Euler characteristic $1$ can be deformed as well and this gives a liftable point in $\CM_0$.

As for (ii), this follows as in the non-relative case. The moduli stack $\mathsf{M}$ of Gieseker-stable 1-dimensional sheaves is a $\BG_m$-gerbe over $\CM$ and, as in \cite{Lieb}, this gerbe is trivial if and only if there exists a twisted line bundle.  In our case, if $\mathsf{F}_T$ is the universal family over $\mathsf{M}$, the determinant $\mathrm{det}R\Gamma(\mathsf{F}_T)$ defines a twisted line bundle because of the condition $\chi(\CF)  =1$.  By pulling back $\mathsf{F}_T$ over the moduli stack via a trivialization of the gerbe, we obtain a universal family $\mathcal{F}_T$ over $\CM$.\footnote{The choice of trivialization affects the universal family by tensoring with a line bundle.}
\end{proof}


\subsection{Specializations}\label{section4.3} Specialization morphisms with respect to perverse filtrations have been studied systematically in \cite{dCM,sp}. Here we provide a concrete description of the specialization morphism in our case as follows.
\source{hitchin-fibrations-abelian-surfaces-pw}

As before, we assume $T = \BP^1$ and $T^\circ = \BP^1 \setminus 0$. Let $f: W \to T$ be a smooth morphism, whose restriction over $T^\circ$ is a trivial product
\[
f^\circ: W^\circ = W_t \times T^\circ \rightarrow T^\circ, \quad \forall~ t\neq 0,
\]
with $W_t$ proper.
Hence by the global invariant cycle theorem \cite[Theorem 1.7.1]{dCM1}, the restriction
\[
\mathrm{res}_t:  H^*(W, \BQ) \to H^*(W_t, \BQ)
\]
is surjective for $t \neq 0$. Let
\[
\alpha_1, \alpha_2 \in H^*(W, \BQ)
\]
be two liftings of a class $\alpha \in H^*(W_t, \BQ)$ with $t\neq 0$. The localization exact sequence implies that
\begin{equation}\label{en52}
\source{hitchin-fibrations-abelian-surfaces-pw}
\alpha_1 - \alpha_2 = {i_0}_* \gamma,
\end{equation}
for some class $\gamma$
 in the Borel--Moore homology of $W_0$.  Here $i_0: W_0 \hookrightarrow W$ is the closed embedding of the central fiber, and we identify the Borel--Moore homology and the cohomology of $W$ in the equation (\ref{en52}). Hence by the excess intersection formula, we have
\[
\mathrm{res}_0(\alpha_1) - \mathrm{res}_0(\alpha_2) = i_0^*{i_0}_* \gamma = c_1(N_{W_0/W})\cap \gamma =0.
\]
We define the specialization morphism
\begin{equation}\label{sp!}
\source{hitchin-fibrations-abelian-surfaces-pw}
\mathrm{sp}^!: H^*(W_t, \BQ) \rightarrow H^*(W_0, \BQ), \quad t\neq 0
\end{equation}
as
\begin{equation}\label{sp1}
\source{hitchin-fibrations-abelian-surfaces-pw}
\mathrm{sp^!}(\alpha) = \mathrm{res}_0(\tilde{\alpha}) \in H^*(W_0, \BQ)
\end{equation}
where $\tilde{\alpha} \in H^*(W, \BQ)$ is any lifting of $\alpha$. The discussion above implies that the class (\ref{sp1}) does not depend on the lifting $\tilde{\alpha}$, and (\ref{sp!}) is a homomorphism of $\BQ$-algebras,
since $\mathrm{res}_t$ is a homomorphism for all $t$.


\begin{example}
\source{hitchin-fibrations-abelian-surfaces-pw}
\notready\label{example1}
Let $\CS \to T$ be the relative surface (\ref{normal_cone}) associated with the embedding $j_C: C \hookrightarrow A$. So we have the specialization morphism
\begin{equation}\label{2333}
\source{hitchin-fibrations-abelian-surfaces-pw}
\mathrm{sp}^!: H^*(A, \BQ) \to H^*(T^*C, \BQ) = H^*(C, \BQ).
\end{equation}
By the definition of $\mathrm{sp}^!$, the morphism (\ref{2333}) is given by the pullback along
\[
T^*C  \hookrightarrow \BP(T^*C \oplus \CO_C) \hookrightarrow \mathrm{Bl}_{C\times 0}(A \times T) \to A\times T \to A.
\]
In particular, we have
\[
\mathrm{sp}^!(\gamma)  = j_C^* \gamma \in H^*(T^*C, \BQ) = H^*(C, \BQ),\quad \forall~ \gamma \in H^*(A, \BQ).
\]
\end{example}

Now we discuss the interaction between $\mathrm{sp}^!$ and perverse filtrations. Let $W$ be as above. We consider a commutative diagram

\begin{equation*}
\text{[Diagram omitted; see source PDF.]}
\end{equation*}
where the morphism $h: W \to V$ is proper of relative dimension $d$. For every $t \in T$, there is a perverse filtration $P_*H^*(W_t, \BQ)$ associated with the morphism $h_t : W_t \to V_t$.

\begin{prop}
\source{hitchin-fibrations-abelian-surfaces-pw}
\notready\label{Prop3.1}
The specialization morphism (\ref{sp!}) preserves the perverse filtrations, \emph{i.e.},
\[
\mathrm{sp}^!\left( P_kH^d(W_t, \BQ ) \right) \subset P_k H^d(W_0, \BQ), \quad \forall k,d.
\]
\end{prop}

\begin{proof}
We show that there exist splittings of the perverse filtrations
\[
H^*(W_t, \BQ) = \bigoplus_{i,d}G_iH^d(W_t, \BQ), \quad \forall~ t \in T
\]
such that
\begin{equation*}
\mathrm{sp}^!\left(G_iH^*(W_t ,\BQ)\right) \subset G_iH^*(W_0, \BQ).
\end{equation*}

We apply the decomposition theorem \cite{BBD} to the morphism $h: W \to V$, and fix an isomorphism
\[
\phi: \mathrm{Rh}_\ast \BQ_W [\mathrm{dim}W - d] \xrightarrow{\simeq} \bigoplus_{i=0}^{2d} \CP_V^i [-i].
\]
Here $\CP_V^i$ are semisimple perverse sheaves on $V$. By the discussion in Section \ref{sec1.3}, the isomorphism $\phi$ induces a splitting of the perverse filtration on $H^*(W, \BQ)$ associated with $h$. By smoothness of $f$ and compatibility of vanishing cycles with $\mathrm{Rh}_*$, we have $\varphi_f(\mathrm{Rh}_\ast \BQ_W) = 0$.  As a result, Corollary $3.1.6$ of \cite{dCM} shows the restriction of $\CP_V^i$ to a closed fiber over $t\in T$ is a shifted perverse sheaf,
\[
\CP_{V,t}^i [-1] = \mathrm{res}_t^*\CP_V^i [-1] \in \mathrm{Perv}(V_t), \quad \forall~ t\in T.
\]
Hence the restriction of the decomposition $\phi$ induces a splitting of the perverse filtration
\[
H^*(W_t, \BQ) = \bigoplus_{i,d}H^d(B, \CP_{B,t}^i)
\]
for every $t\in T$, and we conclude from the definition of $\mathrm{sp}^!$ that
\[
\mathrm{sp}^! (H^d(V, \CP_{V,t}^i)) \subset H^d(V, \CP_{V,0}^i)
\]
This completes the proof.
\end{proof}

Now we consider the family
\[
h_T: \CM \to  \CB
\]
over $T$ constructed in Section \ref{Section 4.2}. Proposition \ref{Prop3.1} implies that the specialization morphism
\[
\mathrm{sp}^!: H^*(\CM_{\beta, A}, \BQ) \to H^*(\CM_{\mathrm{Dol}}, \BQ)
\]
preserves the perverse filtrations. Let $\widetilde{G}_*H^*(\CM_{\beta,A}, \BQ)$ be a splitting given by Theorem \ref{strengthened2.1}. So it is the first Deligne splitting induced by a Lefschetz class
\[
\eta_{A,\beta} \in H^2(\CM_{\beta,A}, \BQ).
\]

\begin{comment}
%associated with $h_t = \pi_\beta: \CM_{A,\beta} \to B$ and $h_0: \CM_{\mathrm{Dol}} \to \Lambda$.

%\begin{prop}
%Let $\hat{G}_*H^*(W_t, \BQ)$ be a splitting of the perverse filtration $P_*H^*(W_t, \BQ)$ for $t\neq 0$. Then we have the decomposition
%\[
%\mathrm{Im}\left\{\mathrm{sp}^!: H^*(W_t, \BQ) \to H^*(W_0,\BQ)\right\}= \bigoplus_{i,d} \mathrm{sp}^!\left(G_iH^d(W_0, \BQ)\right)
%\]
%which splits the restricted perverse filtration $P_*H^*(W_0, \BQ) \cap \mathrm{Im}(\mathrm{sp}^!)$.
%\end{cor}

\begin{proof}
Since $\mathrm{sp}^!$ preserves the perverse filtrations, it suffices to show that
\[
\mathrm{sp}^!\left(G_iH^d(W_0, \BQ)\right) \cap \mathrm{sp}^!\left(G_jH^d(W_0, \BQ)\right) = \{0\}, \quad \forall~ i\neq j.
\]
This follows from Proposition \ref{Prop3.1} that
\[
\mathrm{sp}^!\left(G_iH^d(W_0, \BQ)\right) \cap \mathrm{sp}^!\left(P_{i-1}H^d(W_0, \BQ)\right) = \{0\}. \qedhere
\]
\end{proof}
\end{comment}

\begin{prop}
\source{hitchin-fibrations-abelian-surfaces-pw}
\notready\label{prop4.4}
The class
\[
\eta_0 = \mathrm{sp}^!\left( \eta_{A,\beta}\right) \in H^2(\CM_{\mathrm{Dol}}, \BQ)
\]
is a Lefschetz class with respect to the Hitchin fibration $h: \CM_{\mathrm{Dol}} \to \Lambda$. Assume further that $\widetilde{G}_\ast H^*(\CM_{\mathrm{Dol}}, \BQ)$ is the first Deligne splitting associated with $\eta_0$. Then we have
\begin{equation}\label{prop4.4eq}
\source{hitchin-fibrations-abelian-surfaces-pw}
\mathrm{sp}^!\left( \widetilde{G}_kH^d(\CM_{\beta,A}, \BQ) \right) \subset \widetilde{G}_k H^d(\CM_{\mathrm{Dol}}, \BQ).
\end{equation}
\end{prop}


\begin{proof}
Let $F_{\beta,A} \subset \CM_{\beta,A}$ be a closed fiber of the morphism $\pi_\beta: \CM_{A,\beta} \to B$. We have
\[
 \mathrm{dim}(F_{A,\beta})= \frac{1}{2}\mathrm{dim}(\CM_{A,\beta}) = \tilde{g}.
\]
Since the fiber class
\[
[F_{A,\beta}] \in H^{2\tilde{g}}(\CM_{A,\beta}, \BQ)
\]
lies in $P_0H^{2\tilde{g}}(\CM_{A,\beta}, \BQ)$ by Proposition \ref{prop1.1}, and $\eta_{A,\beta}$ is a Lefschetz class, we obtain that
\[
\eta_{A,\beta}^{\tilde{g}}|_{F_{A,\beta}} = \eta_{A,\beta}^{\tilde{g}} \cup [F_{A,\beta}] \neq 0.
\]
by the hard Lefschetz condition. Hence after specialization, we have
\[
\eta_0^{\tilde{g}}|_{F_0} = \eta_{A,\beta}^{\tilde{g}}|_{F_{A,\beta}} \neq 0
\]
with $F_0$ a closed fiber of $h: \CM_{\mathrm{Dol}} \to \Lambda$. We conclude from Proposition \ref{prop3.2_Hitchin} that $\eta_0$ is a Lefschetz class with respect to $h$. The equation (\ref{prop4.4eq}) then follows from Lemma \ref{comparison1}.
\end{proof}


\begin{rmk}
\source{hitchin-fibrations-abelian-surfaces-pw}
\notready\label{remark1}
Proposition \ref{prop4.4} implies that $\mathrm{sp}^!(G_*H^*(\CM_{A,\beta}, \BQ))$ splits the restricted perverse filtration $P_*H^*(\CM_{\mathrm{Dol}}, \BQ) \cap \mathrm{Im(sp^!)}$, \emph{i.e.},
\[
P_kH^*(\CM_{\mathrm{Dol}}, \BQ) \cap \mathrm{Im(sp^!)} = \bigoplus_{i\leq k} \mathrm{sp}^!\left( \widetilde{G}_iH^*(\CM_{A,\beta}, \BQ)\right).
\]
In general, for an arbitrary splitting $G_*H^*(\CM_{A,\beta}, \BQ)$ of the perverse filtration $P_*H^*(\CM_{A,\beta}, \BQ)$, it may not be true that
\[
\mathrm{sp}^!\left({G}_iH^*(\CM_{A,\beta}, \BQ) \right) \cap \mathrm{sp}^!\left({G}_jH^*(\CM_{A,\beta}, \BQ) \right) = \{0\}, \quad \forall i\neq j.
\]
Hence it is crucial to realize the splitting $\widetilde{G}_*H^*(\CM_{A,\beta}, \BQ)$ as a Deligne splitting associated with a Lefschetz class as in Theorem \ref{strengthened2.1}.
\end{rmk}



\subsection{Normalized classes} Our purpose is to apply the specialization morphism to the tautological classes. We first discuss some properties of the normalized classes introduced in Section 0.3.

Following \cite{HRV,HT1,Shende}, the cohomology
\[
H^*(\CM_{\mathrm{Dol}}, \BQ) = H^*(\CM_B, \BQ)
\]
can be understood by the associated $\mathrm{PGL}_n$-character variety $\tilde{\CM}_B$. More precisely, let
\[
H^*(\CM_B, \BQ) \cong H^*(\tilde{\CM}_B, \BQ) \otimes H^*((\BC^*)^{2g}, \BQ)
\]
be the isomorphism established in \cite[Section 1]{HT1} and \cite[Theorem 2.2.12]{HRV}. Every class $w \in H^i(\tilde{\CM}_B, \BQ)$ can be naturally viewed as a class
\[
w = w\otimes 1 \in H^i(\CM_B, \BQ).
\]
The weights of the tautological classes associated with the universal $\mathrm{PGL}_n$-bundle $\CT$ were calculated in \cite{Shende},
\begin{equation}\label{eqn52}
\source{hitchin-fibrations-abelian-surfaces-pw}
\int_\gamma c_k(\CT) \in {^k\mathrm{Hdg}^{i+2k-2}(\CM_B)},\quad \forall~\gamma \in H^i(C, \BQ).
\end{equation}
The following lemma deduces (\ref{taut_weight}) from (\ref{eqn52}).

\begin{lem}
\source{hitchin-fibrations-abelian-surfaces-pw}
\notready\label{lem3.1}
A twisted universal family $(\CU^\alpha, \theta)$ satisfies
\begin{equation}\label{53}
\source{hitchin-fibrations-abelian-surfaces-pw}
\int_\gamma \mathrm{ch}^\alpha_k(\CU) \in {^k\mathrm{Hdg}^{i+2k-2}(\CM_B)},\quad \forall~\gamma \in H^i(C, \BQ),~~ \forall~ k\geq 0,
\end{equation}
if and only if $\mathrm{ch}^\alpha(\CU)$ is normalized.
\end{lem}

\begin{proof}
We define
\[
{^k\mathrm{Hdg}^{d}(C\times \CM_B)} = \bigoplus_{i+j=d} H^i(C, \BQ) \otimes {^k\mathrm{Hdg}^{j}(\CM_B)}.
\]
Then (\ref{53}) is equivalent to
\[
\mathrm{ch}^\alpha(\CU) \in \bigoplus_k {^k\mathrm{Hdg}^{2k}(C\times \CM_B)}.
\]
By a direct calculation using Chern roots, we have
\[
\mathrm{ch}^\alpha(\CU) = \mathrm{ch}(\CT) \cup \mathrm{exp}\left(\frac{c_1(\CU)}{n}+ \alpha\right).
\]
Note that (\ref{eqn52}) is equivalent to \[
\mathrm{ch}(\CT) \in \bigoplus_k {^k\mathrm{Hdg}^{2k}(C\times \CM_B)}.
\]
Since
\[
{^1\mathrm{Hdg}^{2}(C\times \CM_B)} = H^1(C, \BQ)\otimes H^1(\CM_B ,\BQ),
\]
we obtain that (\ref{53}) holds if and only if
\[
\mathrm{ch}_1^\alpha(\CU)= c_1(\CU) + n\alpha \in H^1(C, \BQ) \otimes H^1(\CM_B, \BQ).
\]
This is equivalent to the condition that the class $\mathrm{ch}(\CU^\alpha)$ is normalized.
\end{proof}

We also give a criterion of the normalized classes $\mathrm{ch}^\alpha(\CU)$ via the perverse filtration parallel to Lemma \ref{lem3.1}.

\begin{lem}
\source{hitchin-fibrations-abelian-surfaces-pw}
\notready\label{lem4.7}
Suppose we have a twisted universal family $(\CU^\alpha, \theta)$ such that, for each $k \geq 0$ and $\gamma \in H^{2i}(C, \BQ)$, the tautological class
\begin{equation*}
\int_\gamma \mathrm{ch}^\alpha_k(\CU) \in H^{2i+2k-2}(\CM_\mathrm{Dol}, \BQ),
\end{equation*}
has perversity $k$. Then $\mathrm{ch}^\alpha(\CU)$ is normalized.
\end{lem}

\begin{proof}
We have
\[
 P_1H^2(\CM_{\mathrm{Dol}}, \BQ) = 0, \quad  P_0H^0(\CM_{\mathrm{Dol}}, \BQ)=H^0(\CM_{\mathrm{Dol}}, \BQ).
\]
For the first vanishing, we can use the decomposition \eqref{112233}; the corresponding vanishing holds $\hat{CM}_{\mathrm{Dol}}$ since its second cohomology is generated by an ample class which cannot have perversity $1$.  The second statement follows similarly.
Since any K\"unneth factor of $\mathrm{ch}_1^\alpha(\CU)$ in $H^*(\CM_{\mathrm{Dol}}, \BQ)$ has perversity $1$, we conclude that
\[
\mathrm{ch}_1^\alpha(\CU) \in  H^1(C, \BQ)\otimes H^1(\CM_\mathrm{Dol} ,\BQ). \qedhere
\]
\end{proof}

\subsection{Specializations and tautological classes}
Assume $j_C: C\hookrightarrow A$ is the closed embedding. Let $\CS \to T$ and $\CM \to T$ be the family of surfaces and the relative moduli space of 1-dimensional sheaves introduced in Section \ref{Section 4.2}. By Lemma \ref{lem4.1}, there is a universal family $\CF_T$ on $\CS\times_T \CM$.

Let
\[
\pi_\CS: \CS\times_T \CM \to \CS,\quad \pi_\CM: \CS\times_T \CM \to \CM
\]
be the projections. The support of $\CF_T$ is the universal spectral curve
\[
\CC \subset \CS\times_T \CM
\]
which is proper over $\CM$. By \cite{BFM}, the sheaf $\CF_T$ defines a localized Chern character
\[
\mathrm{ch}(\CF_T) \in H^{\mathrm{BM}}_*(\CC, \BQ)
\]
with $H^{\mathrm{BM}}_*(-, \BQ)$ the Borel--Moore homology. For any class
\begin{equation}\label{alphatilde}
\source{hitchin-fibrations-abelian-surfaces-pw}
\tilde{\alpha} = \pi_\CS^* \tilde{\alpha}_1 + \pi_\CM^* \tilde{\alpha}_2 \in H^2(\CS \times_T \CM)
\end{equation}
with $\tilde{\alpha}_1\in H^2(\CS, \BQ),~ \tilde{\alpha}_2\in H^2(\CM, \BQ)$, we consider the relative twisted class
\begin{equation}\label{456}
\source{hitchin-fibrations-abelian-surfaces-pw}
\mathrm{ch}^{\tilde{\alpha}}(\CF_T) = \mathrm{exp}(\tilde{\alpha})\cap \mathrm{ch}(\CF_T).
\end{equation}
Since $\mathrm{ch}(\CF_T)$ defines a class in the Borel--Moore homology of $\CC$, the class (\ref{456}) also lies naturally in $ H^{\mathrm{BM}}_*(\CC, \BQ)$, whose degree $k$ part
\[
 \mathrm{ch}_k^{\tilde{\alpha}}(\CF_T) \in H^{\mathrm{BM}}_{2\mathrm{dim}(\CS\times_T\CM)-2k}(\CC, \BQ)
\]
induces the relative tautological classes
\begin{equation}\label{427}
\source{hitchin-fibrations-abelian-surfaces-pw}
\int_{\tilde{\gamma}} \mathrm{ch}_k^{\tilde{\alpha}}(\CF_T) = {\pi_\CC}_* ( \pi_{\CS}^*\tilde{\gamma} \cap  \mathrm{ch}_k^{\tilde{\alpha}}(\CF_T) ) \in H^*(\CM, \BQ), \quad \forall \tilde{\gamma}\in H^*(\CS, \BQ).
\end{equation}
Here $\pi_\CC$ is the proper morphism
\[
\pi_\CC: \CC \hookrightarrow \CS \times_T \CM \to \CM,
\]
and therefore the push-forward
\[
{\pi_\CC}_*: H^\mathrm{BM}_*(\CC, \BQ) \to H^\mathrm{BM}_*(\CM, \BQ) \cong H^*(\CM, \BQ)
\]
is well defined. In particular, (\ref{427}) gives a family of cohomology classes over $T$ whose restriction to every closed fiber is a twisted universal class.

\begin{prop}
\source{hitchin-fibrations-abelian-surfaces-pw}
\notready\label{prop4.8}
Let
\[
\int_{\gamma} \mathrm{ch}_k^\alpha (\CF_\beta) \in H^*(\CM_{A,\beta}, \BQ), \quad  \gamma \in H^*(A, \BQ)
\]
be the classes of Theorem \ref{taut_abelian}, then we have
\[
\mathrm{sp}^!\left( \int_{\gamma} \mathrm{ch}_k^\alpha (\CF_\beta) \right) = c(k-1, j_C^*\gamma).
\]
\end{prop}

\begin{proof}
By the definition of $\mathrm{sp}^!$, we have
\[
\mathrm{sp}^!\left( \mathrm{res}_t\left( \int_{\tilde{\gamma}} \mathrm{ch}_k^{\tilde{\alpha}} (\CF_T) \right)\right) = \mathrm{res}_0\left( \int_{\tilde{\gamma}} \mathrm{ch}_k^{\tilde{\alpha}} (\CF_T) \right)
\]
for any $\tilde{\alpha}$ of the type (\ref{alphatilde}). Hence the class $\mathrm{sp}^!\left( \int_{\gamma} \mathrm{ch}_k^\alpha (\CF_\beta) \right)$ is of the form
\[
\int_{\mathrm{sp}^!(\gamma)} \mathrm{ch}_k^{\alpha_0} (\CF_0),
\]
where $(\CF_0, \alpha_0)$ is a twisted universal family of 1-dimensional Gieseker-stable sheaves on $T^*C \times \CM_{\mathrm{Dol}}$. Example \ref{example1} and a direct calculation using the Grothendieck--Riemann--Roch formula imply that
\begin{equation*}\label{997}
\source{hitchin-fibrations-abelian-surfaces-pw}
\int_{\mathrm{sp}^!(\gamma)} \mathrm{ch}_k^{\alpha_0} (\CF_0) = \int_{j_C^*\gamma} \mathrm{ch}_{k-1}^{\alpha_0}(\CU).
\end{equation*}
Here $\CU = \mathrm{pr}_* \CF_0$ is a universal family on $C\times \CM_{\mathrm{Dol}}$ with
\[
\mathrm{pr}: T^*C\times \CM_{\mathrm{Dol}} \to C\times \CM_{\mathrm{Dol}}
\]
the natural projection.

In conclusion, we obtain that
\[
\mathrm{sp}^!\left( \int_{\gamma} \mathrm{ch}_k^\alpha (\CF_\beta) \right) = \int_{j_C^*\gamma} \mathrm{ch}_{k-1}^{\alpha_0}(\CU).
\]
Moreover, we know from Proposition \ref{prop4.4} that the twisted class \[
\int_{j_C^*\gamma} \mathrm{ch}_{k-1}^{\alpha_0}(\CU),\quad \forall \gamma \in H^*(A, \BQ)
\]
has perversity $k-1$ for any $k \geq 1$. Hence Lemma \ref{lem4.7} implies that
\[
\int_{j_C^*\gamma} \mathrm{ch}_{k-1}^{\alpha_0}(\CU) = c(k-1, j_C^*\gamma). \qedhere
\]
\end{proof}

\subsection{Proofs of Theorems \ref{genus_2}, \ref{P=W_even}, and \ref{P=W_odd}}\label{Section4.6}
\source{hitchin-fibrations-abelian-surfaces-pw}

Since the perverse filtration with respect to the Hitchin fibration $h: \CM_{\mathrm{Dol}} \to \Lambda$ is locally constant when we deform the curve $C$ \cite{dCM}, it suffices to prove all the three theorems for a curve which can be embedded in an abelian surface
\[
j_C: C\hookrightarrow A.
\]
Therefore we apply the specialization morphism
\[
\mathrm{sp}^!: H^*(\CM_{\beta, A}, \BQ) \to H^*(\CM_{\mathrm{Dol}}, \BQ).
\]
introduced in Section \ref{section4.3}, and obtain the following results immediately from Theorem \ref{strengthened2.1}, Proposition \ref{prop4.4}, and Proposition \ref{prop4.8}:
\begin{enumerate}
    \item[(i)] We have
    \begin{equation}\label{333333}
\source{hitchin-fibrations-abelian-surfaces-pw}
    c(\gamma, k) \in \widetilde{G}_kH^*(\CM_{\mathrm{Dol}}, \BQ), \quad \forall \gamma \in \mathrm{Im}({j_C^*}) \subset H^*(C, \BQ).
    \end{equation}
    \item[(ii)] The restriction of the decomposition \[
    H^*(\CM_{\mathrm{Dol}}, \BQ)= \bigoplus_{k,d}\widetilde{G}_kH^d(\CM_{\mathrm{Dol}}, \BQ)
    \]to the subalgebra of $H^*(\CM_{\mathrm{Dol}}, \BQ)$ generated by the classes (\ref{333333}) is multiplicative.
\end{enumerate}

We first prove Theorem \ref{genus_2}. The Abel-Jacobi map embeds a genus 2 curve $C$ into its Jacobian
\[
j_C: C \hookrightarrow \mathrm{Jac}(C) = A.
\]
Hence the restriction morphism $j_C^*$ is surjective, and Theorem \ref{genus_2} follows from (i) and (ii) above.

The proof of Theorem \ref{P=W_even} is similar. For any embedding $j_C$, the image of $j_C^*$ always contains the sub-vector space
\[
H^0(C, \BQ) \oplus H^2(C, \BQ) \subset H^*(C, \BQ).
\]
Hence the subalgebra
\[
R^*(\CM_{\mathrm{Dol}}) \subset H^*(\CM_{\mathrm{Dol}}, \BQ)
\]
is contained in the subalgebra generated by the classes (\ref{333333}), and we again conclude Theorem \ref{P=W_even} by (i) and (ii).

Finally we treat the odd classes
\[
c(\gamma, k) \in H^{2k-1}(\CM_{\mathrm{Dol}}, \BQ), \quad \gamma\in H^1(C, \BQ)
\]
and prove Theorem \ref{P=W_odd}.

When the curve $C$ has genus $\geq 3$, the restriction
\begin{equation}\label{resH1}
\source{hitchin-fibrations-abelian-surfaces-pw}
j_C^*: H^1(A, \BQ) \to H^1(C, \BQ)
\end{equation}
is not surjective. We know from (i) that $c(\gamma, k)$ has perversity $k$ for any $\gamma$ lying in the image of (\ref{resH1}). Since the monodromy group of the moduli space $\CM_g$ of nonsingular genus g curves is the full symplectic group $\mathrm{Sp}_{2g}$ by \cite{A}, the sub-vector space
\[
\mathrm{Im}\left({j_C^*}:H^1(A, \BQ)\to H^1(C,\BQ) \right) \subset H^1(C, \BQ)
\]
generate the total cohomology $H^1(C, \BQ)$ via the action of the monodromy group. We deduce Theorem \ref{P=W_odd} from \cite{dCM}. \qed
